\section{Математическое описание моделей}

\subsection{Латентно-семантический анализ (LSA)}

Математический фундамент LSA строится на принципах линейной алгебры и
статистического анализа частот слов. Процесс начинается с построения
терм-документной матрицы \( X \), где строки - уникальные слова, а столбцы -
фрагменты текста (документы или параграфы). Каждая ячейка матрицы содержит
частоту вхождения слова, которая для повышения точности обычно подвергается
лог-энтропийному взвешиванию (log-entropy weighting). Это позволяет снизить
значимость общеупотребительных слов и выделить специфические термины

Центральным инструментом LSA является сингулярное разложение (SVD), которое
раскладывает исходную матрицу \( X \) на произведение трех матриц:
\[
  X = U \Sigma V^T
\]
\begin{itemize}
  \item \( U \): матрица термов с ортонормированными векторами слов
  \item  \( \Sigma \): диагональная матрица сингулярных значений, упорядоченных
    по убыванию
  \item \( V^T \): матрица документов с ортонормированными векторами
\end{itemize}

LSA строит оптимальное приближение матрицы \( X \) ранга \( k \)
(обычно 200-400 измерений) путем зануления меньших сингулярных
значений. Этот процесс минимизирует норму Фробениуса (L2-норму)
ошибки. Главное ограничение LSA заключается в том, что он работает по
принципу \enquote{мешка слов} (bag-of-words): семантика текста
аппроксимируется как простая линейная сумма векторов входящих в него
слов, полностью игнорируя их порядок и сложные синтаксические связи

Косинусное сходство. После L2-нормализации эмбеддингов
\(\hat{\mathbf{v}}_i = \frac{\mathbf{v}_i}{\|\mathbf{v}_i\|_2}\) мерой
близости двух документов служит косинус угла между ними:

\begin{displaymath}
  \text{sim}(\mathbf{v}_i, \mathbf{v}_j) = \hat{\mathbf{v}}_i^\top
  \hat{\mathbf{v}}_j = \cos\theta_{ij}
\end{displaymath}

\subsection{BERT}

В отличие от линейной аппроксимации LSA, BERT основан на многослойной
архитектуре Transformer, использующей механизмы внимания для динамического
вычисления контекста. Основным математическим компонентом является Scaled
Dot-Product Attention:
\[
  Attention(Q, K, V) = softmax(\frac{QK^T}{\sqrt{d_k}})V
\]

где \( Q, K, V \) - матрицы запросов, ключей и значений, а \( d_k \) -
размерность векторов. Благодаря многоголовному вниманию (Multi-Head Attention)
модель способна параллельно анализировать информацию из различных
репрезентативных подпространств.

Ключевое отличие BERT, обеспечивающее его семантическое качество,
заключается в
нелинейности и контекстуализации:
\begin{itemize}
  \item Позиционное кодирование: В отличие от LSA, BERT вводит
    векторы на основе
    синусоид разной частоты, что позволяет математически учитывать
    порядок слов и
    структуру предложения.

  \item Глубокая двунаправленность: Через задачу Masked LM (MLM)
    модель обучается
    предсказывать замаскированные токены, одновременно учитывая и
    левый, и правый
    контекст во всех слоях. Это позволяет BERT создавать уникальные
    векторы для
    каждого вхождения слова (решая проблему полисемии), тогда как в LSA слово
    всегда имеет один фиксированный вектор.

  \item Нелинейные преобразования (FFN): В каждом слое BERT присутствуют
    полносвязные сети прямого распространения с нелинейными
    функциями активации
    (ReLU или GELU):
    \[
      FFN(x) = \max(0, xW_1 + b_1)W_2 + b_2
    \]
    Именно этот аппарат позволяет модели выходить за рамки простого
    суммирования
    значений слов и улавливать сложные иерархические зависимости, недоступные
    линейным методам.
\end{itemize}
